% 第 63 章骨架：导引框已定，正文待写。写作计划见 docs/PLAN.md“第 63 章”，写法见 docs/STYLE.md。
\chapter{另一条线：明斯基、存量–流量一致模型与 MMT}\label{ch:minsky}

\begin{chapterguide}
\guidequestion{稳定本身为什么会孕育不稳定？}
\begin{guidegoals}
\item 概括明斯基的金融不稳定假说，以及对冲、投机、庞氏三类融资；
\item 构造一个存量—流量一致模型（Godley \& Lavoie 2007），并核对它的账户；
\item 概括现代货币理论的主要命题（国定货币、功能财政、就业保障）及其思想来源；
\item 整理主流对现代货币理论的主要批评，以及现代货币理论学者的回应。
\end{guidegoals}
\guideprereq{\chref{ch:warfinance}\enspace\nameref{ch:warfinance}；\chref{ch:accounts}\enspace\nameref{ch:accounts}；\chref{ch:money}\enspace\nameref{ch:money}}
\end{chapterguide}

\chapterstub
